\documentclass[10pt,a4paper]{amsart}
\usepackage[margin=19mm]{geometry}
\usepackage{amsmath,amssymb,mathtools}
\usepackage[scr=rsfs,cal=euler]{mathalfa}
\usepackage{xfrac}
\usepackage[unicode,hidelinks,bookmarks=false]{hyperref}
\newcommand{\ie}{i.e.\ }
\newcommand{\cf}{cf.\ }
\newcommand{\resp}{respectively\ }
\newcommand{\Subsection}{\subsection}
\providecommand{\nobreakdash}{\nobreak\hbox{-}}
\setlength{\emergencystretch}{2em}
\multlinegap0pt
\usepackage{amssymb}
\newtheorem{lemma}{Lemma}
\newtheorem{corollary}{Corollary}
\newtheorem{theorem}{Theorem}
\newcommand{\diam}{\operatorname{diam}}
\newcommand{\height}{\operatorname{height}}
\newcommand{\spn}{\operatorname{span}}
\title{Sharp order-preserving integer models for short additive equalities}
\author{Enkai Zhang}
\date{Source: arXiv:2609.08915v1 (CC BY 4.0)}
\begin{document}
\maketitle

\noindent\emph{Source and licence.} Sharp order-preserving integer models for short additive equalities. Version arXiv:2609.08915v1 (CC BY 4.0), distributed by arXiv under the Creative Commons Attribution 4.0 International licence, http://creativecommons.org/licenses/by/4.0/, as recorded in the arXiv licence metadata for this exact version and shown on the abstract page https://arxiv.org/abs/2609.08915v1. Full licence notice: \texttt{licenses/arxiv-2609.08915-CC-BY-4.0.txt}.\par\smallskip

\noindent\textit{Editorial scope.} Counted unit: the article's theorem M-thm:lifting (source0.tex lines 1077--1082). The article's complete original proof of it is delivered with it (source0.tex lines 1083--1150, one \texttt{proof} environment of 68 lines). No mathematics is written, repaired or re-derived: the statement and the proof are the article's own lines, in the article's own notation, and no retained line is substituted or re-flowed.  Retained uncounted as the source-local context the proof needs: lines 361--364; lines 380--387; lines 408--408; lines 432--436; lines 452--461.  Nothing was omitted from any retained span, so the package carries no omission record.  No retained line contains a citation, so the wrapper carries no bibliography and no bibliography entry.  No retained line contains a figure environment, an image-inclusion command or drawing code, so no image is delivered and the package declares no asset dependency.  Editorial formatting only, in addition: an AMS \texttt{amsart} preamble that restates the article's own declarations and macros used by the retained lines, namely the environments \texttt{lemma}, \texttt{corollary}, \texttt{theorem} on separate counters, together with the standard packages those lines need.  The retained environments print in this document as Lemma 1, Corollary 1, Theorem 1 in order of appearance, whereas the article prints the same items with the numbering of its own full document.  The complete native source of the article as distributed in the arXiv e-print for this exact version is bundled unchanged as \texttt{sources/209762/source0.tex}.
\begin{lemma}\label{M-lem:mass-minor}
For any square matrix whose rows are deleted rows of balanced relations of
masses $h_1,\ldots,h_s$, $|\det|\le\prod_{i=1}^sh_i$.
\end{lemma}

\begin{corollary}[one-dimensional kernels]\label{M-cor:max-rank}
Suppose $r=m-2$, so that $d=1$. Then $L$ is a rational line, and $x$ is a
positive multiple of the primitive integer generator $w$ of $L$ with positive
last coordinate. Consequently $T$ is affinely equivalent to the primitive
integer alphabet $\{0,w_1,\ldots,w_{m-1}\}$, which is an order-preserving
model of $T$ preserving all balanced relations, and its diameter satisfies
$w_{m-1}\le H\le q^{m-2}$.
\end{corollary}

\section{Flags in the ordered relation kernel}\label{M-sec:flag}

\begin{lemma}[chamber extension]\label{M-lem:chamber-extension}
Let $G$ be a face of $K_0$ of dimension at least two, let $F$ be a facet of
$G$, and let $C_F$ be a chamber of $F$. Then there is a chamber $C$ of $G$
with $C\cap F=C_F$; in particular $C_F$ is a facet of $C$.
\end{lemma}

\begin{theorem}[Flag bound]\label{M-thmD}\label{M-thm:flag-model}
For $m\ge2$ and $q\ge2$, the alphabet has an order-preserving $q$-equivalent
primitive integer model of diameter at most
\[
 H\sum_{j=0}^{d-1}q^j.
\]
Here the basis is chosen from the vanishing integer $q$-relation rows.
The masses in $H$ belong to their full balanced vectors. The space $V$
contains all vanishing $q$-relations.
\end{theorem}

\begin{theorem}[face-lifting recurrence]\label{M-thm:lifting}
For $m\ge3$ and $q\ge2$,
\begin{equation}\label{M-eq:recurrence}
 H_m(q)\le q^{m-2}+H_{m-1}(q).
\end{equation}
\end{theorem}

\begin{proof}
We first model the coincident letters on an ordering facet, then lift
that model to a compatible chamber and check its order, relations and diameter.
Given an $m$-letter real alphabet, let $V$ be its complete $q$-relation span,
and let $K=K_0$ be the weak-order cone in its translated real kernel $L$, as
in Section~\ref{M-sec:flag}. It has full relative dimension $d$ in $L$ and
contains a strictly ordered point.

If $d=1$, a primitive ray generator gives diameter at most $q^{m-2}$ directly
by the balanced-row cofactor bound (Corollary~\ref{M-cor:max-rank}), so
\eqref{M-eq:recurrence} follows. Suppose $d\ge2$ and take an original ordering
facet $F$ of $K$, that is, a facet of the form $K\cap\{x_1=0\}$ or
$K\cap\{x_i=x_{i+1}\}$. Its nonzero relative interior has a fixed pattern of
coincident adjacent letters, hence $k$ distinct value blocks, for some
$2\le k\le m-1$: there are at least two blocks, because a single common value
with the first translated coordinate fixed at zero would give only the zero
vector, and there are at most $m-1$ blocks, because a proper ordering facet
has at least one active gap equality. Several equalities may be forced
simultaneously on $F$, so $k$ may be smaller than $m-1$.

\smallskip\noindent\emph{Model on the face.}
Choose $p$ in that relative interior avoiding every $q$-hyperplane that does
not contain $\spn F$. This is possible because there are finitely many
proper hyperplanes. Let $0=\pi_1<\cdots<\pi_k$ be the distinct values among
$0,p_1,\ldots,p_{m-1}$ and $\beta(i)$ the block of letter $i$, so that
$p_i=\pi_{\beta(i)}$. Choose an order-preserving $q$-model
$0=\pi'_1<\cdots<\pi'_k$ of $(\pi_1,\ldots,\pi_k)$ of diameter at most
$H_k(q)$ and set $u_i=\pi'_{\beta(i)}$. The resulting integer vector $u$
then has $\height(u)=\diam\le H_k(q)\le H_{m-1}(q)$.

It lies in $F$. Indeed every original $q$-relation collapses, by summing its
coefficients within each block, to a balanced relation on the $k$ blocks of
mass no larger than $q$, since summing within blocks only permits
cancellation; its value at the repeated vector equals the value of the
collapsed relation at the block values. The collapsed relation vanishes at
the block values of $p$, so the block model preserves it, and the original
relation vanishes on $u$; as these relations span $V$, $u\in L$. All the
required equalities between adjacent letters are preserved by repetition, and
all other ordering gaps are positive. Likewise every $q$-form not identically
zero on $F$ remains nonzero at $u$: it was nonzero at $p$ and collapses to a
relation of mass at most $q$ on the value blocks, whose nonvanishing the
block model preserves. Thus $u$ is in the relative interior of a full
short-hyperplane chamber $C_F$ on $F$, even when $F$ has multiple forced
equalities among the original $m$ coordinates. Its nonzero comparison signs
need not match those of $p$; what follows uses the chamber that actually
contains $u$.

\smallskip\noindent\emph{Lift to a chamber.}
Lemma~\ref{M-lem:chamber-extension} supplies an adjacent $d$-dimensional
chamber $C$ inside $K$ with $C\cap F=C_F$. Choose a primitive integer extreme
ray $v$ of $C$ outside $\spn F$; one exists because $C$ spans $L$. It is
determined by $m-2$ independent original, short-relation and ordering rows,
each of mass at most $q$: an extreme ray of the $d$-dimensional pointed
cone $C$ lies on $d-1$ independent active hyperplanes of its defining list,
by the perturbation argument in the proof of Theorem~\ref{M-thm:flag-model},
and these together with the $r$ rows spanning $V$ are $m-2$ independent
rows. So its height is at most $q^{m-2}$ by Lemma~\ref{M-lem:mass-minor}.

\smallskip\noindent\emph{Order, relations and diameter.}
The vector $u+v$ is in the relative interior of $C$. A supporting functional
of $C$ that vanished on $u$ must vanish on the entire span of $F$, and $v$
makes it strict; every other facet functional is already strict on $u$.
Equivalently, $u$ supplies all directions within $F$ and $v$ leaves that
facet. Hence $u+v$ is strictly ordered and lies on no $q$-hyperplane outside
$V$, while all original relations vanish on it. Gcd division gives the
desired primitive model, of height at most $H_{m-1}(q)+q^{m-2}$. This proves
\eqref{M-eq:recurrence}.
\end{proof}

\end{document}
